\section{Run-disjoint OpenTelemetry experiment}
\label{sec:v11-experiment}

\subsection{System and collection protocol}

We use an unmodified checkout of the OpenTelemetry Demo.  The experiment records
the repository commit, working-tree status, Docker version, Docker Compose
version, and the Compose files used for execution.  The application is started
once and the load generator is restarted before each collection window.  A run
is one warm-up interval followed by one fixed trace-collection interval.
Queries are restricted by the corresponding microsecond timestamps, so traces
from different runs are not mixed.

The collector first queries the Jaeger service endpoint and resolves the
Checkout service name from the returned service list.  This avoids assuming a
fixed spelling across Demo releases.  The raw Jaeger response for each run is
retained without modification.

\subsection{Checkout synchronization projection}

Only traces containing explicit Checkout CLIENT to downstream SERVER parent
links are retained.  We do not infer missing cross-service links.  For a normal
checkout containing $n$ product items, the projection is
\[
  w_n=p^n\#c^n s\#(pc)^n s.
\]
The first component is the Product Catalog SERVER lifeline, the second is the
Currency SERVER lifeline, and the third is the Checkout CLIENT schedule.
The final Currency call in caller order is designated as the shipping
conversion; preceding Currency calls are product-price conversions.

\subsection{Route-preserving offline mutations}

Mutations are constructed from normal traces after collection.  No application
source code is changed.  Every mutation keeps the caller schedule component
fixed and alters only downstream SERVER order or multiplicity.  We consider:
(i) inversion of the final product-conversion and shipping-conversion SERVER
order; (ii) deletion or duplication of a Product Catalog SERVER occurrence;
(iii) deletion or duplication of a product-price Currency SERVER occurrence;
and (iv) duplication of the shipping Currency SERVER occurrence.  A deletion is
called route preserving only if another occurrence of the same service edge
remains.

\subsection{Strict run-disjoint evaluation}

For each random seed, complete collection runs are partitioned into training,
validation, and test sets.  No trace or mutant derived from one run crosses a
split boundary.  The grammar anchors are the normal projections for cart sizes
one and two in the training runs.  Observer selection uses the smallest
available unseen normal cart size and one phase-inversion mutant whose cart size
is represented in the positive anchors.  The test set is not consulted during
observer selection.

We report normal acceptance separately for cart sizes represented in the
positive anchors and for unseen cart sizes.  Mutation recall is similarly split
into seen-size and unseen-size subsets.  This prevents a model from receiving
credit for rejecting a mutant merely because it rejects every string of that
length.

\subsection{Compared models}

The main comparison contains: the selected fan-out-three MCFG, the same learner
under the trivial observation, the full direct product of all candidate
observers, an exact-template model, and 2-gram and 3-gram anomaly models.  For
each MCFG we report the total rule count, unit-rule count, and number of observed
types.  Seed-level results and aggregate mean, standard deviation, minimum, and
maximum are retained.

\subsection{Controlled pipeline validation}

Before the official-Demo run, we validate the complete pipeline on seven
independent Jaeger-format fixture runs.  Across five run-split seeds, the greedy
selector chooses only the shipping-phase count modulo two observer.  The
selected grammar accepts all unseen normal cart sizes and rejects all six
mutation families.  The trivial-observation grammar retains normal
extrapolation but rejects only half of the seen-size phase inversions, whereas
the full observer product and exact-template model reject all unseen normal cart
sizes.  These values validate the implementation and are not reported as
results on the official OpenTelemetry Demo.
